\section*{Permutations, support and fresh quantification (publisher sections 2.1--2.3)}
We fix disjoint countably infinite sets: the set of atoms $\mathbb A=\{a,b,c,\ldots\}$ and unknowns $\mathbb V=\{X,Y,Z,\ldots\}$.
A (finite) permutation on $\mathbb A$ is a bijection $\pi:\mathbb A\to\mathbb A$ whose set of moved atoms $\operatorname{mv}(\pi)=\{a\in\mathbb A:\pi(a)\ne a\}$ is finite. Write $\pi\circ\pi'$, $\mathsf{id}$ and $\pi^{-1}$ for functional composition, the identity, and the inverse of permutations, respectively.
The group of all finite permutations of atoms is denoted $G(\mathbb A)$.
Every finite permutation $\pi\in G(\mathbb A)$ can be uniquely decomposed (up to reordering) into pairwise disjoint cycles as $\pi=\gamma_1\circ\cdots\circ\gamma_n$, with $\operatorname{mv}(\gamma_i)\cap\operatorname{mv}(\gamma_j)=\varnothing$ for $i\ne j$.
This is called the canonical form of $\pi$.
\paragraph{Assumption 2.1.1.} Permutations are always in a canonical form.
A nominal set is a pair $(\mathcal X,\cdot)$ consisting of a set $\mathcal X$ and an action $\cdot:G(\mathbb A)\times\mathcal X\to\mathcal X$ such that every element $x\in\mathcal X$ has finite support, i.e., there exists a finite $B\subseteq\mathbb A$ with $\pi\cdot x=x$ for all $\pi\in G(\mathbb A)$ such that $\pi(a)=a$ for every $a\in B$.
The least such $B$ is denoted $\operatorname{supp}(x)$. An atom $a$ is fresh for $x$ (written $a\#x$) if $a\notin\operatorname{supp}(x)$.
Sometimes, when there is no ambiguity, we refer to a nominal set $(\mathcal X,\cdot)$ by its carrier set $\mathcal X$ and leave the action $\cdot$ implicit.
Throughout this section, let $\mathcal X$ denote a nominal set. We consider formulas of the form $\phi(c,\bar x)$, where $c$ is an atom and $\bar x=x_1,\ldots,x_n$ is a finite list of parameters with each $x_i\in\mathcal X$.
Typically, $\phi$ is a predicate in higher-order logic [22] or ZFA set theory [14,16].
\paragraph{Definition 2.3.1.} The $\mathsf N$-quantifier is formally defined as:
\[
\newq c.\phi(c,\bar x)\ \Longleftrightarrow\ \{c\in\mathbb A:\phi(c,\bar x)\}\text{ is cofinite}.
\]
A cofinite subset of a set $M$ is a subset $B$ whose complement in $M$ is a finite set.
In cases where multiple fresh atoms are introduced at once, we consider nested $\mathsf N$-quantification of the form $\newq c_1\cdots\newq c_n.\phi(c_1,\ldots,c_n,\bar x)$, abbreviated as $\newq\bar c.\phi(\bar c,\bar x)$.
For a finite set of atoms $A$ and a finite list of elements $\bar x=x_1,\ldots,x_k$ of $\mathcal X$, write $A\#\bar x$ as abbreviation for $A\cap\operatorname{supp}(\bar x)=\varnothing$, where $\operatorname{supp}(\bar x)=\bigcup_{j=1}^k\operatorname{supp}(x_j)$.
If $\bar a=a_1,\ldots,a_n$ is a list of atoms, we write $\bar a\#x$ as shorthand for $\{\bar a\}\#x$, where $\{\bar a\}=\{a_1,\ldots,a_n\}$.
A fundamental connection between the $\mathsf N$-quantifier and freshness is given by the following equivalence [16]:
\[
a\#x\ \Longleftrightarrow\ \newq c.(a\ c)\cdot x=x.
\]
We define a generalised fixed-point identity as an assertion of the form $\newq\bar c.\pi\cdot x=x$, where $\pi$ is a permutation that may contain atoms bound by $\mathsf N$.
\paragraph{Definition 2.3.2 ($\bar c$-cycles and $\bar c$-permutations).}
Let $\pi$ be a permutation written as a product of disjoint cycles and let $\bar c$ be a list of atoms.
\begin{enumerate}
\item A cycle $\eta$ of $\pi$ is called a $\bar c$-cycle if $\operatorname{mv}(\eta)\cap\{\bar c\}\ne\varnothing$, and a $\neg\bar c$-cycle otherwise.
\item A permutation $\pi$ is a $\bar c$-permutation iff all its cycles are $\bar c$-cycles, and a $\neg\bar c$-permutation iff all its cycles are $\neg\bar c$-cycles.
\end{enumerate}
Write $\pi=\pi_{\bar c}\circ\pi_{\neg\bar c}$ where $\pi_{\bar c}$ is the product of all $\bar c$-cycles of $\pi$, and $\pi_{\neg\bar c}$ is the product of all $\neg\bar c$-cycles of $\pi$.
For a set $S\subseteq G(\mathbb A)$, we write $\langle S\rangle$ for the subgroup generated by $S$, that is, the smallest subgroup of $G(\mathbb A)$ containing $S$.
In particular, if $S=\{\rho_1,\ldots,\rho_n\}$ is a finite set of permutations, we write $\langle\rho_1,\ldots,\rho_n\rangle$ for $\langle S\rangle$.

\section*{Complete author supporting proofs (publisher appendix A)}
\paragraph{Lemma 2.3.4.} Let $x\in\mathcal X$ and $\pi,\eta_1,\ldots,\eta_n$ be permutations written as compositions of disjoint cycles.
\begin{enumerate}
\item If $\pi\cdot x=x$, $\eta$ is a cycle of $\pi$, and $\operatorname{mv}(\eta)-\operatorname{supp}(x)\ne\varnothing$, then $\operatorname{mv}(\eta)\#\operatorname{supp}(x)$.
\item If $\eta_i\cdot x=x$, for all $i=1,\ldots,n$, and $\rho\in\langle\eta_1,\ldots,\eta_n\rangle$, then $\rho\cdot x=x$.
\end{enumerate}
\begin{proof}
\begin{enumerate}
\item Write $\pi=\mu_1\circ\cdots\circ\mu_n$ with disjoint cycles $\mu_1,\ldots,\mu_n$, and let $\eta=\mu_j$ for some $1\le j\le n$. Choose $a\in\operatorname{mv}(\eta)\setminus\operatorname{supp}(x)$.
Suppose, by contradiction, there exists $b\in\operatorname{mv}(\eta)\cap\operatorname{supp}(x)$.
Since the cycles are disjoint, $\mu_i(b)=b$ for all $i\ne j$. Hence $\pi(b)=\eta(b)$.
Since $b\in\operatorname{supp}(x)$ and $\pi\cdot x=x$, it follows that $\pi(b)\in\operatorname{supp}(\pi\cdot x)=\operatorname{supp}(x)$, and therefore $\eta(b)\in\operatorname{supp}(x)$.
Because $a,b\in\operatorname{mv}(\eta)$, there exists $l>0$ such that $\eta^l(b)=a$.
Applying $\pi$ repeatedly to $\eta(b)\in\operatorname{supp}(x)$ and using that $\pi\cdot x=x$ yields $\eta^l(b)\in\operatorname{supp}(x)$, hence $a\in\operatorname{supp}(x)$, contradicting the choice of $a$.
\item Assume that $\eta_i\cdot x=x$ for every $i=1,\ldots,n$, and let $\rho\in\langle\eta_1,\ldots,\eta_n\rangle$.
By definition of the subgroup generated by $\eta_1,\ldots,\eta_n$, the permutation $\rho$ can be written as a finite composition $\rho=\zeta_1\circ\cdots\circ\zeta_k$, where each $\zeta_r$ belongs to $\{\eta_1^{\pm1},\ldots,\eta_n^{\pm1}\}$.
We first note that each inverse $\eta_i^{-1}$ also fixes $x$.
Indeed, from $\eta_i\cdot x=x$, applying $\eta_i^{-1}$ to both sides gives $\eta_i^{-1}\cdot(\eta_i\cdot x)=\eta_i^{-1}\cdot x$, and therefore $x=\eta_i^{-1}\cdot x$.
So every $\zeta_r$ in the above decomposition fixes $x$.
We now prove that $\rho$ fixes $x$ by repeated application of these equalities:
\[
\rho\cdot x=(\zeta_1\circ\cdots\circ\zeta_k)\cdot x=\zeta_1\cdot(\zeta_2\cdot(\cdots(\zeta_k\cdot x)\cdots)).
\]
Since $\zeta_k\cdot x=x$, the innermost term simplifies to $x$; then $\zeta_{k-1}\cdot x=x$, and continuing in this way we obtain $\rho\cdot x=x$.
\end{enumerate}
\end{proof}
A key property of the $\mathsf N$-quantifier is the Some/Any principle [16, Prop.\ 4.10]:
\paragraph{Proposition A.1 (Some/Any Principle).} The following are equivalent:
\[
\newq c.\phi(c,\bar x)\ \Longleftrightarrow\ \forall c.(c\#\bar x\Rightarrow\phi(c,\bar x))\ \Longleftrightarrow\ \exists c.(c\#\bar x\land\phi(c,\bar x)).
\]
For multiple atoms, nested $\mathsf N$-quantifiers introduce distinctness conditions. We avoid writing these explicitly by fixing the following convention.
\paragraph{Notation A.2.} Let $\bar c=c_1,\ldots,c_n$ be a list of pairwise distinct atoms.
We abbreviate $\forall c_1\cdots\forall c_n.(\bigwedge_{i<j}c_i\ne c_j\Rightarrow\phi)$ as $\forall\bar c.\phi$, and similarly for $\exists\bar c.\phi$.
\paragraph{Lemma A.3 (Generalised Some/Any Principle).} Let $\bar c$ be a list of atoms. Then:
\[
\newq\bar c.\phi(\bar c,\bar x)\ \Longleftrightarrow\ \forall\bar c.(\bar c\#\bar x\Rightarrow\phi(\bar c,\bar x))\ \Longleftrightarrow\ \exists\bar c.(\bar c\#\bar x\land\phi(\bar c,\bar x)).
\]
\begin{proof}
By induction on the length of $\bar c$. The base case is (4).
For the inductive step, write $\bar c'=\bar c,c_{n+1}$ and define $\psi(\bar c,\bar x):=\newq c_{n+1}.\phi(\bar c,c_{n+1},\bar x)$.
Then $\newq\bar c'.\phi\equiv\newq\bar c.\psi$ and by the induction hypothesis and (4) this is equivalent to $\forall\bar c\ \forall c_{n+1}.(\bar c\#\bar x\land c_{n+1}\#\bar x,\bar c\Rightarrow\phi)$.
By Notation A.2, this is exactly $\forall\bar c'.(\bar c'\#\bar x\Rightarrow\phi)$, as required. The existential case is analogous.
\end{proof}

\paragraph{Theorem 2.3.5 (Generalised Fix-Point).} Let $x\in\mathcal X$ and let $\pi$ be a permutation.
\begin{enumerate}
\item If $\pi$ is a $\bar c$-permutation and $A=\operatorname{mv}(\pi)\setminus\{\bar c\}$, then $\newq\bar c.\pi\cdot x=x$ iff $A\#x$.
\item $\newq\bar c.\pi\cdot x=x$ iff $(\newq\bar c.\pi_{\bar c}\cdot x=x)\land(\pi_{\neg\bar c}\cdot x=x)$.
Equivalently, if $A=\operatorname{mv}(\pi_{\bar c})\setminus\{\bar c\}$, then $\newq\bar c.\pi\cdot x=x$ iff $A\#x\land(\pi_{\neg\bar c}\cdot x=x)$.
\end{enumerate}

\begin{proof}
\begin{enumerate}
\item $(\Rightarrow)$ Assume $\newq\bar c.\pi\cdot x=x$.
By Lemma A.3, for some (equivalently, any) choice of $\bar c$ such that $\bar c\#x$, we have $\pi\cdot x=x$.
Let $\eta$ be a cycle of $\pi$. Since $\pi$ is a $\bar c$-permutation, $\eta$ contains some atom $a\in\bar c$.
As $a\#x$, Lemma 1 yields $\operatorname{mv}(\eta)\#x$. Hence $\operatorname{mv}(\pi)\#x$, and therefore $A\#x$.
$(\Leftarrow)$ Assume $A\#x$. Let $\bar c$ be such that $\bar c\#x$.
Since $\pi$ is a $\bar c$-permutation, every cycle of $\pi$ contains an atom from $\bar c$ and all remaining atoms lie in $A$.
Thus every cycle is disjoint from $\operatorname{supp}(x)$, so $\pi\cdot x=x$.
By Lemma A.3, it follows that $\newq\bar c.\pi\cdot x=x$.
\item If $\pi_{\bar c}=\mathsf{id}$ or $\pi_{\neg\bar c}=\mathsf{id}$, the result is immediate. Assume both are non-identity.
We first show that $\newq\bar c.\pi\cdot x=x$ iff $\newq\bar c.(\pi_{\bar c}\cdot x=x\land\pi_{\neg\bar c}\cdot x=x)$.
By Lemma A.3, both sides are equivalent to universal statements over $\bar c$ with the freshness condition $\bar c\#x,\operatorname{mv}(\pi)\setminus\{\bar c\}$.
Hence it suffices to show that, for such $\bar c$, one has $\pi\cdot x=x$ iff both $\pi_{\bar c}\cdot x=x$ and $\pi_{\neg\bar c}\cdot x=x$.
Assume $\pi\cdot x=x$. Let $\eta$ be a cycle of $\pi_{\bar c}$. Since $\pi_{\bar c}$ is a $\bar c$-permutation, $\eta$ contains some atom $a\in\bar c$.
As $a\#x$, Lemma 1 yields $\operatorname{mv}(\eta)\#x$. Since this holds for every cycle $\eta$ of $\pi_{\bar c}$, and the cycles are disjoint, we obtain $\operatorname{mv}(\pi_{\bar c})\#x$, and therefore $\pi_{\bar c}\cdot x=x$ (equivalently, $\pi_{\bar c}^{-1}\cdot x=x$).
From $\pi=\pi_{\bar c}\circ\pi_{\neg\bar c}$ and $\pi\cdot x=x$, we have $\pi_{\bar c}\cdot(\pi_{\neg\bar c}\cdot x)=x$.
Applying $\pi_{\bar c}^{-1}$ to both sides and using that $\pi_{\bar c}^{-1}\cdot x=x$, it follows that $\pi_{\neg\bar c}\cdot x=x$.
Conversely, if both $\pi_{\bar c}$ and $\pi_{\neg\bar c}$ fix $x$, then $\pi\cdot x=(\pi_{\bar c}\circ\pi_{\neg\bar c})\cdot x=\pi_{\bar c}\cdot(\pi_{\neg\bar c}\cdot x)=x$.
From the fact that $\mathsf N$ distributes over conjunction, we obtain $\newq\bar c.(\pi_{\bar c}\cdot x=x\land\pi_{\neg\bar c}\cdot x=x)$ is equivalent to $\newq\bar c.(\pi_{\bar c}\cdot x=x)\land\newq\bar c.(\pi_{\neg\bar c}\cdot x=x)$.
Since $\pi_{\neg\bar c}$ does not involve atoms from $\bar c$, the formula $\pi_{\neg\bar c}\cdot x=x$ is independent of $\bar c$, and hence, $\newq\bar c.(\pi_{\neg\bar c}\cdot x=x)$ is equivalent to $\pi_{\neg\bar c}\cdot x=x$.
Therefore, $\newq\bar c.\pi\cdot x=x$ iff $\newq\bar c.(\pi_{\bar c}\cdot x=x)\land(\pi_{\neg\bar c}\cdot x=x)$, and the final equivalence follows from item (1), applied to $\pi_{\bar c}$ and $B$.
\end{enumerate}
\end{proof}

